\subsection{Two finite class congruences for the entire complex norm system}

\begin{statement}{E.88}{Proved}{All outside-S divisor lifts, convergent product coordinates, and the complex positional cube condition}
Use the actual complex fields $L_d$ and roots of E.72,
including $L_0=K=\mathbb Q(\zeta)$. Put
\[
 \mathcal X=\varprojlim_{d,N}\widehat{L_d^\times}_3,\qquad
 \mathcal X_S=\varprojlim_{d,N}
       \bigl(\mathcal O_{L_d}[1/3,1/17]^\times\otimes\mathbb Z_3\bigr).
 \tag{E88.1}
\]
The inverse limits use the actual field norms. Let $\mathcal P$
be the rational primes $p\ne3,17$ with
$v_3(\operatorname{ord}_p17)=1$. For each such $p$, choose the
unique residue value of $\zeta$ for which $17\zeta$ has order
prime to three, and the unique prime-to-three roots
$r_{d,p}^{3^d}=17\zeta$ in $\mathbb F_p^\times$. Write
$\mathfrak P_{d,p}=(p,\beta_d-r_{d,p})$.
These are precisely the coherent prime chains of E72.3.
The notation $\mathcal P$ is the rational-prime indexing
of the set $\mathcal P_1$ there.

Here is a finite-field and degree-six algorithm for the
entire class map. Compute the class of
$\mathfrak P_{1,p}$ in the certified class group
$\operatorname{Cl}(L_1)=C_3^2$ of E.57, obtaining $(u_p,v_p)$.
In the coherent basis of E.86 put
\[
 \rho(p)=(c_p,d_p)=(2u_p+2v_p,v_p)\pmod3,\qquad
             c_p,d_p\in\{0,1,2\}.
 \tag{E88.2}
\]
This is an evaluation in one fixed number field, with
its retained integral basis, ideal class generators and
unconditionally certified class map. In particular
$\rho(7)=(2,1)$ and $\rho(73)=(0,1)$.
For
\[
 \mathscr R(\mathcal P)
 =\{(t_p)\in\prod_{p\in\mathcal P}\mathbb Z_3:
                         t_p\longrightarrow0\pmod{3^m}
                         \text{ for every }m\},
\]
define the continuous surjection
\[
 \rho:\mathscr R(\mathcal P)\longrightarrow C_3^2,\qquad
       (t_p)\longmapsto
         \left(\sum c_pt_p,\sum d_pt_p\right)\pmod3,
 \qquad \mathcal K=\ker\rho .
 \tag{E88.3}
\]
The sums modulo three have finite support. The exact
global divisor sequence is
\[
 0\longrightarrow\mathcal X_S
 \longrightarrow\mathcal X
 \xrightarrow{\ \operatorname{div}_{\notin S}\ }
       \mathscr R(\mathcal P)
 \xrightarrow{\ \rho\ }C_3^2
 \longrightarrow0 .
 \tag{E88.4}
\]
Thus the two congruences in (E88.3) are necessary and
sufficient for a coherent multiplicative lift at all
levels, allowing valuations at $3,17$.

There are explicit topological coordinates on their
kernel. If $e_p$ denotes its divisor coefficient vector,
put
\[
 E_7=3e_7,\qquad E_{73}=3e_{73},\qquad
 E_p=e_p+c_pe_7-(c_p+d_p)e_{73}
                       \quad(p\ne7,73).
 \tag{E88.5}
\]
The following map and inverse are topological
isomorphisms:
\[
 \begin{split}
 \mathbb Z_3^2\times\mathscr R(\mathcal P\setminus\{7,73\})
       &\longrightarrow\mathcal K,\\
 (z_7,z_{73},(r_p))
       &\longmapsto z_7E_7+z_{73}E_{73}+\sum r_pE_p,\\
 r_p&=t_p\quad(p\ne7,73),\\
 z_7&=\frac{t_7-\sum c_pt_p}{3},\qquad
 z_{73}=\frac{t_{73}+\sum(c_p+d_p)t_p}{3}.
 \end{split}
 \tag{E88.6}
\]
The sums in the inverse run over $p\ne7,73$ and converge
in $\mathbb Z_3$. Their numerators are divisible by three
by exactly (E88.3).

Choose, for each vector $E_i$ in (E88.5), one actual
coherent lift $x^{(i)}\in\mathcal X$ supplied by (E88.4).
Remove its common valuation at seventeen by dividing by
the corresponding power of the coherent tuple
$\beta=(\beta_d)_{d\ge0}$; all chosen lifts will have
valuation zero there. Then
\[
 \begin{split}
 s:\mathcal K&\longrightarrow\mathcal X,\qquad
 s(t)=\bigl(x^{(7)}\bigr)^{z_7}
       \bigl(x^{(73)}\bigr)^{z_{73}}
           \prod_{p\ne7,73}\bigl(x^{(p)}\bigr)^{r_p},\\
 \mathcal X_S\times\mathcal K&\xrightarrow{\ \sim\ }\mathcal X,
                   \qquad (u,t)\longmapsto u\,s(t),\\
 x&\longmapsto
   \bigl(x\,s(\operatorname{div}_{\notin S}x)^{-1},
                       \operatorname{div}_{\notin S}x\bigr).
 \end{split}
 \tag{E88.7}
\]
The product is convergent in every finite field and
every power quotient, hence in the displayed inverse
limit. Both maps are continuous. Changing the choices
multiplies $s$ by a continuous homomorphism
$\mathcal K\to\mathcal U_\infty$, where
$\mathcal U_\infty$ is the coherent $3$-unit group
of E.69. Conversely every such homomorphism gives
another normalized splitting.

In particular E.69 and E72.10 give the full compact
kernel and its literal coordinates:
\[
 \mathcal X_S=\beta^{\mathbb Z_3}\mathcal U_\infty
 \simeq \mathbb Z_3^2
                \times\prod_{d\ge0}\mathbb Z_3^{\,2\cdot3^d}.
 \tag{E88.8}
\]
The two factors in (E88.7), together with (E88.6)
and E69.6, compute the entire complex multiplicative
norm system as a topological group with specified
forward and inverse maps. The restricted sequence
factor is given its multiplication-by-$3^m$
prodiscrete topology; it is not declared compact.

The rational base norm has a precise additional
coordinate. Define
\[
 \Theta(t)=v_{1-\zeta}\bigl(s(t)_0\bigr)\in\mathbb Z_3 .
 \tag{E88.9}
\]
This is a continuous linear map. A normalized
splitting change modifies it by $3^h$ times a
continuous linear map to $\mathbb Z_3$, with the
actual finite integer $h$ of E.69. Thus
$\Theta\bmod3^h$ is independent of the choices.
There is a fully evaluated reduction modulo three.
Let $(\ell_p,m_p)\in C_9\oplus C_3$ be the
ordinary class of $\mathfrak P_{2,p}$ in E86.8,
with $0\le\ell_p<9$. It is computed from the
same lower prime $\mathfrak P_{1,p}$ by the
retained genus Artin map; its first coordinate
modulo three is $c_p$. Then
\[
 \Theta(t)\equiv
       -\frac{\sum_{p\in\mathcal P}\ell_p t_p}{3}
          \pmod3\qquad(t\in\mathcal K).
 \tag{E88.10}
\]
The numerator belongs to $3\mathbb Z_3$, and
changing an integer representative of $\ell_p$
does not change the residue.

Here is the exact transfer to global $H$-section
classes. In E.46 fix $A=3\kappa_3(17)$ and take
pure positional component $1$. The entire complex
component locus satisfying the total-norm cube
condition consists of
\[
 \begin{gathered}
 x=\beta^{\,3b}\,u(z,\mathbf a)\,s(3r),\qquad
 b,z\in\mathbb Z_3,\quad
 \mathbf a\in\prod_{d\ge0}\mathbb Z_3^{\,2\cdot3^d},\\
 r\in\mathscr R(\mathcal P),\qquad
                         \sum c_pr_p=0\pmod3 .
 \end{gathered}
 \tag{E88.11}
\]
Here $u(z,\mathbf a)$ is the actual recursive
coherent $3$-unit map E69.6. Formula (E88.7)
gives the inverse and proves uniqueness of these
coordinates. In these coordinates the scaled
toric map E46.7 is
\[
 \left(
 \kappa_3(17),\
 2b\,\kappa_3(17)
 +\left(3^{h-1}z+\frac{\Theta(3r)}3\right)\kappa_3(3)
 +\sum_{p\in\mathcal P}r_p\kappa_3(p)
 \right).
 \tag{E88.12}
\]
The quotient $\Theta(3r)/3$ is integral by
(E88.10)--(E88.11). The subset of $r$ in
(E88.11) itself has the explicit restricted
basis
$3e_7$ and $e_p-\widetilde{2c_p}e_7$ for
$p\ne7$, where $\widetilde{2c_p}\in\{0,1,2\}$
is the indicated residue. Thus this entire
locus has coordinates
$\mathbb Z_3^3\times\mathscr R(\mathcal P\setminus\{7\})
\times\prod_{d\ge0}\mathbb Z_3^{2\cdot3^d}$.
These are the section classes in the actual
target $H$ of the quotient maps E.38 and
the Kummer--Shapiro comparison E.46.

Two literal norm-coherent examples fix useful
arithmetic origins:
\[
 x_d^+=\beta_d+1,\qquad x_d^-=\beta_d-1,\qquad
 N_{K/\mathbb Q}x_0^+=273=3\cdot7\cdot13,\quad
 N_{K/\mathbb Q}x_0^-=307.
 \tag{E88.13}
\]
Their outside-S divisors are $e_7+e_{13}$ and
$e_{307}$, and their valuations at three are
$1$ and $0$. Their cubes therefore give actual
global $H$ sections with second toric coordinates
$\kappa_3(273)$ and $\kappa_3(307)$, respectively.
The certificate evaluates all $79$ eligible
primes at most $2000$, the complete finite pivot
matrix and its inverse, and the phase (E88.10).

The remaining integer $h$ is retained exactly
as in E.69. The values of the choice-independent
map $\Theta\bmod3^h$ beyond its displayed
modulo-three evaluation are not claimed
numerically evaluated here. The previous
all-level principal-class equations away S
are replaced by the two explicitly evaluated
conditions (E88.3); no unknown higher
S-class matrix remains. The manuscript
anchor is LM~\S51.16, fols.~733v--736v,
scans~172--178, with the actual maps
E.46, E.69, E.72 and E.86.
\end{statement}

\begin{proof}
\emph{The actual class maps at every level.}
For $p\in\mathcal P$, its indicated residue
root and its cube define the exact prime
norm
$N_{L_{d+1}/L_d}\mathfrak P_{d+1,p}
=\mathfrak P_{d,p}$, with residue degree
one. E.86 identifies every S-class norm
with the identity in the coherent
$(a_d,b_d)$ basis. Therefore the S-class
of this prime has the same two coordinates
at every level. At level one the columns
of $(a_1,b_1)$ in the certified class
basis are $(2,0),(2,1)$. The inverse
matrix modulo three is
$\left(\begin{smallmatrix}2&2\\0&1\end{smallmatrix}\right)$,
which proves (E88.2). It proves the
two necessary class congruences at
every level, not just at the bottom.

There is a complete finite algorithm
for these coordinates for any eligible
prime. Set $q=(p-1)/3^{v_3(p-1)}$.
For the two primitive cubic roots in
$\mathbb F_p$, exactly one makes
$17\zeta$ have order prime to three.
Exponentiation by the inverse of
$3^d$ modulo $q$ gives $r_{d,p}$.
The ideal $(p,\beta_1-r_{1,p})$ is
then evaluated by the certified
degree-six class map. The code retains
the maximal-order ideal interpretation
and checks its norm $p$. The official
algorithm returns the coordinates in
the retained class-generator ideals
\cite[entries \texttt{bnfcertify},
\texttt{bnfisprincipal}, \texttt{bnrisprincipal}]
{E_PARICertifiedUnits}.
At the next level the Frobenius in
the lower genus field of E.86 is
computed on this same prime. It lies
in the kernel of the actual cubic
restriction character, because the
specified residue is a cube. Since
that genus field is the full Hilbert
three-class field of $L_2$, E86.8
gives the entire ordinary class.
Norm compatibility gives
$\ell_p\bmod3=c_p$, $m_p=d_p$.

\emph{Completion and sufficiency.}
For one number field $F$, the ordinary
outside-S divisor sequence is
\[
 0\longrightarrow\mathcal O_{F,S}^\times
 \longrightarrow F^\times
 \longrightarrow\operatorname{Div}_{\notin S}(F)
 \longrightarrow\operatorname{Cl}_S(F)
 \longrightarrow0.
 \tag{E88.14}
\]
Its image of principal divisors is
a subgroup of a free abelian group,
hence is free. The first short exact
sequence therefore splits as abelian
groups. For the second short exact
sequence the class quotient is finite.
More explicitly, let $D'$ be the
preimage of the prime-to-three
subgroup of the class group in
the divisor group. Then
$D'/\operatorname{Prin}$ has
prime-to-three order, so the
two subgroups have the same
three-adic completion. Choose $3^e$
annihilating the three-primary
class group. It gives the
cofinal inclusions
$3^{m+e}\operatorname{Div}\subset
3^mD'\subset3^m\operatorname{Div}$.
Thus the
completed sequence remains exact:
\[
 0\longrightarrow\widehat{\mathcal O_{F,S}^\times}_3
 \longrightarrow\widehat{F^\times}_3
 \longrightarrow\widehat{\operatorname{Div}_{\notin S}(F)}_3
 \longrightarrow\operatorname{Cl}_S(F)\{3\}
 \longrightarrow0 .
 \tag{E88.15}
\]
The first term is compact, because
the S-unit group is finitely generated.
Every nonempty fiber of the completed
divisor map is a translate of that
compact group. In particular it is
compact even though the completed
divisor group is prodiscrete with
infinite quotients.

E72.4, with zero coefficients on S,
identifies the outside-S divisor
norm limit with $\mathscr R(\mathcal P)$.
For $t\in\mathcal K$, its divisor
at every level has zero S-class,
by the norm-identical class maps
already proved. Each completed
principal-generator fiber is
therefore nonempty and compact.
Norm maps carry a fiber into its
predecessor. Any finite collection
of compatibility equations has
a solution: choose an element
in the highest involved fiber
and take its successive actual
norms. The product of all fibers
is compact, and these are closed
conditions with the finite
intersection property. Hence
their entire intersection is
nonempty. This supplies one
coherent lift with all the
prescribed outside-S divisors.
Necessity follows from (E88.15).
The kernel is exactly
$\mathcal X_S$, component by
component. Surjectivity of
$\rho$ follows from the two
actual prime columns $(2,1)$,
$(0,1)$ at $7,73$. This proves
every kernel and image in
(E88.4).

\emph{Restricted coordinates and convergent splittings.}
Substitute the two pivot columns
into (E88.5). Every displayed
vector has zero S-class.
For a restricted sequence $t$,
the two sums in (E88.6) converge:
their coefficients are bounded
integers and $t_p$ tends to zero
modulo every power of three.
The two class equations say
exactly that both resulting
numerators are divisible by
three. Reading all nonpivot
coordinates proves the formula
and inverse. Reduction modulo
$3^m$ and, for the inverse,
modulo $3^{m+1}$ proves their
continuity. This proves (E88.6)
in its actual prodiscrete topology.

Choose one lift of each basis
vector using the preceding
compact-fiber argument. The
seventeen-valuation is the same
at every level because that
prime has residue degree one;
dividing by the appropriate
power of $\beta$ removes it
compatibly. For each fixed
level $d$ and quotient by
$3^m$th powers, all but
finitely many factors in
(E88.7) equal one: their
exponents $r_p$ belong to
$3^m\mathbb Z_3$. Thus the
products converge, independently
of their ordering. Norms commute
with the finite products and
with these limits. The product
therefore belongs to $\mathcal X$
and has outside-S divisor $t$.
This proves the homomorphism
and section assertions.

Continuity is equally explicit.
A basic open condition in
$\mathcal X$ imposes a power
quotient at finitely many field
levels. Exponents divisible by
a sufficiently large power of
three send every factor to one
in all these quotients. Thus
the linear product map is
continuous in the uniform
power topology on the restricted
sequence module, and in the
usual topology on the two
finite coordinates. The inverse
in (E88.7) is a composition
of continuous divisor,
section and multiplication
maps. This proves the
topological isomorphism without
a compactness claim for its
restricted factor.

Two choices of a normalized
lift of a basis vector differ
by an actual coherent $3$-unit.
Their products define a
continuous homomorphism
$\mathcal K\to\mathcal U_\infty$
by the same convergence
argument. Conversely any
continuous homomorphism of
these modules is determined
by its values on the displayed
basis: finitely supported
coordinate vectors are dense.
Multiplying the chosen lifts
by those values gives exactly
the corresponding changed
splitting. This proves the
full choice ambiguity. Finally
remove the common seventeen
valuation from an arbitrary
element of $\mathcal X_S$;
the remainder is precisely
the coherent $3$-unit group.
E69.6 supplies its actual
product coordinates and their
inverse, proving (E88.8).

\emph{The base phase and the positional condition.}
Valuation is continuous and
commutes with the degree-one
norms at three, so (E88.9)
is continuous and linear.
E69.3 says that the base
projection of every coherent
$3$-unit is a power
$\lambda^{3^h z}$,
$\lambda=1-\zeta$.
This proves the stated
change of $\Theta$. Conversely
the chosen E69.6 section of
this base projection lifts
every continuous linear
$\mathcal K\to\mathbb Z_3$
to a possible splitting change.
Thus the ambiguity is exactly
the one displayed.

At level two the full divisor
of $s(t)_2$ is
$\Theta(t)\mathfrak q_{3,2}
+\sum t_p\mathfrak P_{2,p}$,
with zero seventeen coefficient.
Its ordinary class must vanish.
E.86 gives
$[\mathfrak q_{3,2}]=(3,0)$
and the coordinates $(\ell_p,m_p)$.
The resulting first coordinate
equation is
$3\Theta(t)+\sum\ell_pt_p=0$
in $C_9$. The second is already
the S-class condition. This
proves (E88.10), including
the integrality of its
division and its independence
of coordinate representatives.

For $x=\beta^{\,b}u(z,\mathbf a)s(t)$,
the actual rational base norm is
\[
 N_{K/\mathbb Q}x_0
   =17^{2b}\,3^{\,3^h z+\Theta(t)}
                    \prod_{p\in\mathcal P}p^{t_p}
       \quad\text{in }\widehat{\mathbb Q^\times}_3 .
 \tag{E88.16}
\]
Its divisor follows from the
displayed valuations and
residue degrees; the possible
$\mu_3$ factor in $K$ has
rational norm one. The
pro-three completion of the
rational unit group is zero,
so these prime coordinates
determine the entire rational
element. The product converges
in that completion by the
restricted-support condition.

This element is a cube exactly
when $b\in3\mathbb Z_3$,
each $t_p\in3\mathbb Z_3$,
and $\Theta(t)\in3\mathbb Z_3$,
since $h\ge1$. Write $t=3r$.
It automatically belongs to
$\mathcal K$. Equation
(E88.10) gives
$\Theta(3r)=-\sum c_pr_p$
modulo three. This proves
both directions of (E88.11).
The first-coordinate cocycle
is still the specified
$A=3\kappa_3(17)$ in E.46;
its actual scaled toric map
divides the rational norm
coordinate by three. Applying
that map to (E88.16) proves
(E88.12), with all factors
and divisibilities typed.
The one-row pivot basis for
the condition on $r$ follows
from $c_7=2$, whose inverse
in $\mathbb F_3$ is $2$.
The same restricted-product
argument proves its forward
and inverse topology.

For the two examples,
$N_{L_{d+1}/L_d}(\beta_{d+1}\pm1)
=\beta_d\pm1$ follows directly
from the odd cubic polynomial.
Their rational base norms
are $17^2-17+1=273$ and
$17^2+17+1=307$. The unique
degree-one coherent branches
give the stated outside-S
divisors, and $\beta_d+1$
is the uniformizer at three
while $\beta_d-1$ is a unit
there. These literal maps
prove (E88.13) and the
two asserted section classes.
The finite certificate checks
the same phase and class
equations, all inverse pivot
coordinates, and the norm
identities independently.
\end{proof}
